\documentclass[12pt, a4paper]{article}

\usepackage[utf8]{inputenc}
\usepackage[T1]{fontenc}
\usepackage[brazilian]{babel}
\usepackage{geometry}
\usepackage{enumitem}
\usepackage{titlesec}
\usepackage{parskip}
\usepackage{hyperref}
\usepackage{xcolor}
\usepackage{tikz}
\usetikzlibrary{arrows.meta, positioning, shapes.geometric}

\geometry{margin=2.5cm}

\definecolor{primary}{HTML}{2563EB}
\definecolor{accent}{HTML}{059669}
\definecolor{muted}{HTML}{6B7280}
\definecolor{danger}{HTML}{DC2626}

\hypersetup{
    colorlinks=true,
    linkcolor=primary,
    urlcolor=primary
}

\titleformat{\section}{\Large\bfseries\color{primary}}{}{0em}{}
\titleformat{\subsection}{\large\bfseries}{}{0em}{}

\title{\Huge\bfseries AgentOS\\[0.3em]\Large\color{muted} O que é esse projeto?}
\author{}
\date{}

\begin{document}

\maketitle
\thispagestyle{empty}

% ─────────────────────────────────────────────
\section{O que são agentes de IA?}

Antes de falar do AgentOS, vale entender o que são \textbf{agentes de IA} --- porque são diferentes do ChatGPT que a maioria das pessoas conhece.

Um chatbot como o ChatGPT é uma conversa: você pergunta, ele responde. Ele não faz nada no seu computador. Já um \textbf{agente de IA} é um programa que pode \textbf{agir} --- ele lê arquivos, escreve documentos, navega na internet, executa comandos e toma decisões sozinho para cumprir uma tarefa. É a diferença entre perguntar ``como eu faço um bolo?'' e ter alguém que vai na cozinha e faz o bolo pra você.

Alguns exemplos reais que já existem hoje:

\begin{itemize}[leftmargin=1.5em]
    \item \textbf{Claude Code} (da Anthropic) --- um agente que lê o código de um projeto, entende a estrutura, escreve código novo, roda testes e corrige erros. Um programador pede ``adicione login nesse sistema'' e o agente faz tudo sozinho.

    \item \textbf{Codex} (da OpenAI) --- similar ao Claude Code. Recebe tarefas de programação e executa de forma autônoma.

    \item \textbf{Devin} --- um agente que tenta replicar o trabalho completo de um desenvolvedor, do planejamento à implementação.
\end{itemize}

Esses agentes já são impressionantes sozinhos. Mas eles criam dois problemas sérios.

% ─────────────────────────────────────────────
\section{Problema 1: Agentes trabalham isolados}

Esses agentes trabalham \textbf{sozinhos}. Cada um opera de forma isolada, sem saber que outros existem.

Se você quer dois agentes fazendo coisas diferentes no mesmo projeto --- por exemplo, um pesquisando e outro implementando baseado nessa pesquisa --- você precisa abrir dois terminais, copiar informações entre eles manualmente, e torcer pra que um não atrapalhe o outro.

Não existe uma forma de dizer: ``Agente~A, faça a pesquisa. Quando terminar, passe o resultado pro Agente~B implementar. Mas antes, me deixa revisar o que o A produziu.''

É como ter uma empresa onde cada funcionário trabalha numa sala isolada, sem conseguir falar com os outros, e o chefe precisa ir de sala em sala levando papéis de um pro outro.

% ─────────────────────────────────────────────
\section{Problema 2: Ninguém controla o que os agentes fazem}

Este é o problema maior --- e o que torna o AgentOS diferente de todas as outras ferramentas no mercado.

Hoje, quando uma empresa deixa seus funcionários usarem agentes de IA, ela não tem como responder perguntas básicas:

\begin{itemize}[leftmargin=1.5em]
    \item \textbf{Quantos agentes estão rodando agora?} Não sei.
    \item \textbf{O que eles estão acessando?} Não sei.
    \item \textbf{Quanto estamos gastando com isso?} Preciso checar manualmente em cada conta de API.
    \item \textbf{Algum agente acessou dados confidenciais?} Impossível saber.
    \item \textbf{Se um auditor perguntar o que nossos agentes fizeram no último mês, temos como mostrar?} Não.
\end{itemize}

Isso se chama \textbf{shadow AI} --- agentes de IA operando dentro de uma organização sem controle, sem visibilidade e sem registro. É o equivalente a dar acesso irrestrito ao computador da empresa para dezenas de assistentes que ninguém supervisiona.

Para uma pessoa usando um agente no computador pessoal, isso pode não parecer grave. Mas para empresas --- especialmente em setores regulamentados como finanças, saúde e direito --- é um risco real. Auditores cobram registros. Reguladores exigem conformidade. Um agente que acessa dados errados ou toma uma decisão não supervisionada pode gerar consequências sérias.

\textbf{Nenhuma ferramenta no mercado resolve esse problema.} As alternativas existentes (CrewAI, AutoGen, LangGraph) focam em fazer agentes conversarem entre si. Nenhuma se preocupa com governança: controle de acesso, trilha de auditoria, limites de orçamento, aprovações obrigatórias, segurança.

% ─────────────────────────────────────────────
\section{A Solução}

O AgentOS é uma \textbf{plataforma de governança e orquestração de agentes de IA}. Pense nele como a combinação de um gerente de projeto com um departamento de compliance --- automatizado:

\begin{itemize}[leftmargin=1.5em]
    \item \textbf{Coordena múltiplos agentes} trabalhando juntos em um fluxo de trabalho definido
    \item \textbf{Controla a ordem de execução} --- quem faz o quê, e em que sequência
    \item \textbf{Permite revisão humana} --- pontos de aprovação onde você revisa o trabalho antes de continuar
    \item \textbf{Controla custos com limites rígidos} --- define orçamentos e o sistema para automaticamente se o limite for atingido. Sem surpresas na fatura
    \item \textbf{Registra absolutamente tudo} --- cada ação, cada arquivo acessado, cada decisão fica salva em um registro imutável que serve como trilha de auditoria
    \item \textbf{Controla permissões} --- cada agente só pode acessar o que foi explicitamente permitido. Um agente de pesquisa não acessa o banco de dados financeiro. Um agente de código não acessa dados de clientes
    \item \textbf{Funciona com qualquer agente} --- não é preso a um fornecedor. Coordena Claude Code, Codex, ou qualquer outro agente de IA
\end{itemize}

A diferença fundamental: outras ferramentas fazem agentes trabalharem juntos. O AgentOS faz agentes trabalharem juntos \textbf{de forma segura, controlada e auditável}.

% ─────────────────────────────────────────────
\section{Como Funciona --- Exemplo Prático}

Imagine que você gerencia investimentos e quer analisar se vale a pena investir em uma empresa. Hoje, você pediria pra diferentes pessoas da sua equipe fazerem partes dessa análise. Com o AgentOS, agentes de IA fazem esse trabalho --- e você continua no controle.

\vspace{1em}
\begin{center}
\begin{tikzpicture}[
    node distance=0.8cm,
    agent/.style={rectangle, rounded corners=6pt, draw=primary, fill=primary!8, text width=3cm, align=center, minimum height=1cm, font=\small\bfseries},
    gate/.style={diamond, draw=accent, fill=accent!10, text width=1.6cm, align=center, aspect=2.2, font=\scriptsize\bfseries, inner sep=1pt},
    arr/.style={-{Stealth[length=6pt]}, thick, color=muted}
]
    \node[agent] (a1) {Analista de\\Mercado};
    \node[gate, right=of a1] (g1) {Sua\\Aprovação};
    \node[agent, right=of g1] (a2) {Analista\\Financeiro};
    \node[gate, right=of a2] (g2) {Sua\\Aprovação};
    \node[agent, right=of g2] (a3) {Relatório\\Final};

    \draw[arr] (a1) -- (g1);
    \draw[arr] (g1) -- (a2);
    \draw[arr] (a2) -- (g2);
    \draw[arr] (g2) -- (a3);
\end{tikzpicture}
\end{center}
\vspace{1em}

\begin{enumerate}[leftmargin=1.5em]
    \item O \textbf{Analista de Mercado} (um agente de IA) pesquisa notícias, tendências do setor, concorrentes e regulamentações. Produz um resumo estruturado com os pontos principais.

    \item O fluxo \textbf{pausa} e mostra esse resumo pra você. Você lê, e decide:
    \begin{itemize}
        \item \textbf{Aprovar} --- ``A pesquisa está boa, pode continuar''
        \item \textbf{Rejeitar com feedback} --- ``Faltou analisar o mercado asiático'' (o agente recebe essa instrução)
    \end{itemize}

    \item O \textbf{Analista Financeiro} (outro agente) recebe a pesquisa aprovada e analisa balanços, receitas, margens, dívidas. Produz uma avaliação financeira com indicadores e riscos.

    \item Pausa de novo pra você revisar os números.

    \item O \textbf{Gerador de Relatório} compila tudo em um documento final com recomendação de investimento.
\end{enumerate}

\textbf{Durante todo esse processo, o AgentOS:}
\begin{itemize}[leftmargin=1.5em]
    \item Registrou cada ação de cada agente no log de eventos
    \item Garantiu que o Analista de Mercado só acessou fontes de notícias (não dados internos)
    \item Garantiu que o Analista Financeiro só acessou os balanços permitidos
    \item Controlou quanto cada agente gastou em chamadas de API
    \item Gerou uma trilha de auditoria completa que mostra exatamente o que aconteceu, quando, e por quê
\end{itemize}

Se um auditor perguntar ``como essa recomendação de investimento foi gerada?'', você tem a resposta completa.

\subsection{O que muda em relação a fazer manualmente?}

\begin{itemize}[leftmargin=1.5em]
    \item \textbf{Velocidade} --- Os agentes trabalham em minutos o que levaria horas ou dias de trabalho humano
    \item \textbf{Consistência} --- O fluxo sempre segue os mesmos passos, nada é esquecido
    \item \textbf{Controle total} --- Você aprova cada etapa --- nada avança sem sua autorização
    \item \textbf{Custo previsível} --- Limites de orçamento rígidos, sem surpresas
    \item \textbf{Auditabilidade} --- Registro completo de tudo que aconteceu --- quem fez o quê, quando, com que dados, e qual foi o resultado
\end{itemize}

% ─────────────────────────────────────────────
\section{Conceitos Principais}

\begin{description}[style=nextline, leftmargin=1.5em]
    \item[\textnormal{\textbf{Fluxo de trabalho (Workflow)}}] Uma sequência de tarefas conectadas. Você define quem faz o quê, em que ordem, e onde precisa de aprovação humana.

    \item[\textnormal{\textbf{Agentes}}] Programas de IA que executam tarefas específicas. O AgentOS coordena agentes de diferentes fornecedores (Anthropic, OpenAI, e outros) através de uma interface única.

    \item[\textnormal{\textbf{Gates (Portões de aprovação)}}] Pontos de controle humano no fluxo. Podem ser aprovações (``Está bom? Sim/Não'') ou pedidos de informação (``Qual empresa analisar?'').

    \item[\textnormal{\textbf{Log de eventos}}] Um registro completo e imutável de tudo que aconteceu. Funciona como uma trilha de auditoria --- nada é perdido, nada é alterado. Serve tanto para monitoramento quanto para conformidade regulatória.

    \item[\textnormal{\textbf{Orçamento}}] Limites configuráveis por agente, por equipe e por fluxo de trabalho. O sistema para automaticamente se o limite for atingido --- sem exceções.

    \item[\textnormal{\textbf{Permissões}}] Cada agente recebe apenas o acesso mínimo necessário para sua tarefa. Um agente de pesquisa pode ler notícias mas não pode acessar o banco de dados interno. Essas restrições são aplicadas pelo sistema, não por instrução ao agente.
\end{description}

% ─────────────────────────────────────────────
\section{Por que isso importa?}

Agentes de IA estão se tornando cada vez mais comuns nas empresas. O mercado está se movendo rápido:

\begin{itemize}[leftmargin=1.5em]
    \item O GitHub adicionou agentes de IA (Claude e Codex) diretamente na plataforma de desenvolvimento
    \item Empresas como Microsoft, Google e Anthropic estão investindo bilhões em agentes autônomos
    \item Equipes de engenharia já usam múltiplos agentes no dia a dia
\end{itemize}

Mas ninguém construiu a \textbf{camada de governança} --- o sistema que garante que esses agentes operem de forma segura, controlada e auditável dentro de uma organização. É como se todo mundo estivesse construindo carros cada vez mais rápidos, mas ninguém tivesse construído as estradas, os semáforos e as leis de trânsito.

O AgentOS é essa infraestrutura. Ele não compete com os agentes --- ele os torna seguros de usar.

% ─────────────────────────────────────────────
\section{Pra Quem é?}

Na primeira versão, o AgentOS é voltado para \textbf{equipes de engenharia de software} que já usam agentes de código (Claude Code, Codex) e querem coordená-los de forma segura. A interface é via linha de comando e arquivos de configuração.

Mas a arquitetura é genérica --- o mesmo motor que coordena agentes de código pode coordenar agentes de qualquer área:

\begin{itemize}[leftmargin=1.5em]
    \item \textbf{Finanças} --- Agentes coletam dados de mercado, analisam balanços e geram relatórios de risco --- com aprovação humana entre cada etapa
    \item \textbf{Pesquisa científica} --- Agentes fazem revisão de literatura, projetam experimentos e analisam resultados
    \item \textbf{Jurídico} --- Agentes pesquisam jurisprudência, redigem pareceres e verificam conformidade
    \item \textbf{Marketing} --- Agentes analisam concorrentes, geram estratégias de conteúdo e criam materiais
\end{itemize}

O que muda entre esses casos é o que cada agente faz. O motor de orquestração, governança, aprovações e auditoria é o mesmo.

% ─────────────────────────────────────────────
\section{O Modelo de Negócio}

O AgentOS não cobra pelas chamadas de IA dos agentes --- essas são pagas diretamente ao fornecedor (Anthropic, OpenAI, etc.). O AgentOS cobra pela plataforma de governança e orquestração:

\begin{itemize}[leftmargin=1.5em]
    \item \textbf{Gratuito} para desenvolvedores individuais usando localmente
    \item \textbf{Plano de equipe} para times que precisam de coordenação e monitoramento
    \item \textbf{Plano empresarial} para organizações que precisam de conformidade, auditoria, controle de permissões e implantação interna
\end{itemize}

Isso significa que o custo do AgentOS é previsível e independente do volume de uso dos agentes --- similar a como se paga por ferramentas de monitoramento ou segurança, não por processamento.

% ─────────────────────────────────────────────
\section{O Diferencial Competitivo}

Existem outras ferramentas de orquestração de agentes (CrewAI, AutoGen, LangGraph). O AgentOS se diferencia em quatro pontos:

\begin{enumerate}[leftmargin=1.5em]
    \item \textbf{Governança como fundação, não como funcionalidade adicional.} O AgentOS foi projetado desde o início para registrar tudo, controlar permissões e garantir auditabilidade. Outras ferramentas adicionam isso depois --- e de forma incompleta.

    \item \textbf{Funciona com agentes reais e autônomos.} Outras ferramentas coordenam cadeias de prompts simples. O AgentOS coordena agentes que operam de verdade --- Claude Code, Codex --- que escrevem código, acessam arquivos e executam comandos.

    \item \textbf{Não depende de um fornecedor.} O AgentOS coordena agentes de qualquer fornecedor. Uma empresa que usa Claude e Codex juntos não pode depender de uma solução da Anthropic ou da OpenAI --- precisa de uma plataforma neutra.

    \item \textbf{Pacotes de conformidade por setor.} No futuro, o AgentOS oferecerá configurações pré-montadas para setores regulamentados (finanças, saúde, jurídico) --- com regras de auditoria, controle de dados e relatórios prontos para inspeção.
\end{enumerate}

% ─────────────────────────────────────────────
\section{Estado Atual}

O projeto está em início de desenvolvimento da versão~1. O foco dos primeiros meses é:

\begin{enumerate}[leftmargin=1.5em]
    \item Construir o motor principal --- log de eventos, executor de fluxos de trabalho, controle de orçamento
    \item Integrar com o Claude Code como primeiro agente suportado
    \item Implementar portões de aprovação e revisão humana
    \item Demonstrar o primeiro fluxo completo: dois agentes colaborando com aprovação humana entre eles
\end{enumerate}

A meta é ter uma versão funcional em 6 meses, focada em equipes de engenharia de software como primeiro mercado.

\end{document}
